\section{Four equivalent coordinate systems}
\label{sec:coordinates}
%======================================================================

Different calculations become simple in different coordinates.  The point is not to multiply notation, but to move between product, factor-gap, norm, and null coordinates without losing information.

Let
\begin{equation}
 w=x+iy
 =cq+i\frac{q^2-c^2}{2},
 \qquad
 \rho:=|w|=\sqrt{x^2+y^2}.
 \label{eq:w-rho}
\end{equation}
Then
\begin{equation}
 \rho=\frac{c^2+q^2}{2}.
 \label{eq:rho-cq}
\end{equation}
The most useful coordinate charts are collected in \cref{tab:coordinates}.

\begin{table}[htbp]
\centering
\caption{Equivalent coordinates for a canonical squared factor point.}
\label{tab:coordinates}
\renewcommand{\arraystretch}{1.35}
\begin{tabular}{>{\raggedright\arraybackslash}p{0.21\textwidth} >{\raggedright\arraybackslash}p{0.31\textwidth} >{\raggedright\arraybackslash}p{0.34\textwidth}}
\toprule
Coordinate system & Definition & Exact inverse information\\
\midrule
Factor coordinates & $(c,q)$ & $m=cq$, $q=\Pplus(m)$\\
Fermat coordinates & $a=(c+q)/2$, $b=(q-c)/2$ & $c=a-b$, $q=a+b$\\
Squared Cartesian & $x=cq$, $y=(q^2-c^2)/2$ & $w=x+iy=(a+ib)^2$\\
Radial-null & $U=\rho-y$, $V=\rho+y$ & $U=c^2$, $V=q^2$, $x^2=UV$\\
Square-root & $a^2=(\rho+x)/2$, $b^2=(\rho-x)/2$ & sign of $b$ is sign of $y$\\
\bottomrule
\end{tabular}
\end{table}

The central identities are
\begin{equation}
 \boxed{
 U:=\rho-y=c^2,
 \qquad
 V:=\rho+y=q^2.
 }
 \label{eq:UV}
\end{equation}
Consequently,
\begin{equation}
 x=\sqrt{UV}=cq,
 \qquad
 y=\frac{V-U}{2}.
 \label{eq:xy-UV}
\end{equation}

These null coordinates linearize the arithmetic cutoffs.  They also make the meaning of $2ab$ exact:
\begin{equation}
 \boxed{Y=2ab=\frac{V-U}{2},}
 \label{eq:Y-null}
\end{equation}
so the imaginary coordinate is half the separation between the squared canonical factors.

%======================================================================
